\documentclass[11pt]{article}
% =========================================================
%  學生講義 共用前言（以 XeLaTeX 編譯）
%  需要套件：xeCJK, tcolorbox（其餘多在 BasicTeX 內）
% =========================================================
\usepackage[a4paper,margin=2.3cm]{geometry}
\usepackage{amsmath,amssymb}
\usepackage{xcolor}
\usepackage{graphicx}

% ---- 中文字型（macOS 系統字型）----
\usepackage{xeCJK}
\setCJKmainfont{Songti TC}      % 內文：宋體
\setCJKsansfont{Heiti TC}       % 標題/強調：黑體
\xeCJKsetup{AutoFakeBold=2.2, AutoFakeSlant=0.2}

% ---- 版面 ----
\setlength{\parindent}{0pt}
\setlength{\parskip}{0.55em}
\linespread{1.15}
\renewcommand{\baselinestretch}{1.15}

% ---- 顏色 ----
\definecolor{cBlue}{HTML}{1f6feb}
\definecolor{cGray}{HTML}{6b7280}
\definecolor{cGrayBg}{HTML}{f3f4f6}
\definecolor{cPurp}{HTML}{7a3ff2}
\definecolor{cPurpBg}{HTML}{f5f1ff}

% ---- 標註框 ----
\usepackage[most]{tcolorbox}

% 就地補充新概念（灰底）：\begin{補充框}{標題}
\newtcolorbox{補充框}[1]{
  enhanced, breakable, arc=2pt, boxrule=0.6pt,
  colback=cGrayBg, colframe=cGray,
  left=9pt,right=9pt,top=7pt,bottom=7pt,
  before upper={\textbf{\sffamily #1}\par\vspace{3pt}},
}

% 延伸 / 開放問題（紫色左邊條）：\begin{延伸框}{標題}
\newtcolorbox{延伸框}[1]{
  enhanced, breakable, arc=2pt, boxrule=0pt,
  colback=cPurpBg, colframe=cPurpBg,
  borderline west={3pt}{0pt}{cPurp},
  left=10pt,right=9pt,top=7pt,bottom=7pt,
  before upper={\textbf{\sffamily 延伸閱讀｜#1}\par\vspace{3pt}},
}

% ---- 圖：置中、底下小字說明 ----
% \fig{檔名}{寬度}{說明}
\newcommand{\fig}[3]{%
  \begin{center}
  \includegraphics[width=#2\linewidth]{#1}\\[2pt]
  {\small\color{cGray} #3}
  \end{center}}

% ---- 章節標題用黑體 ----
\newcommand{\h}[1]{\sffamily #1}


\begin{document}

\begin{center}
{\sffamily\Huge\bfseries 必物-4 電與磁的統一}\\[3pt]
{\sffamily\large 普通高中 物理（必修）・學生講義}
\end{center}
\vspace{0.6em}

本章先認識「電」：電荷如何透過\textbf{電場}互相施力，並用\textbf{庫侖定律}把這個力寫成公式。接著認識「磁」：電流會產生磁場（電生磁），變動的磁場又會產生電場（磁生電）。最後是十九世紀的綜合——\textbf{馬克士威（Maxwell）}把電與磁整理成一組方程式，並從中推導出一種以光速行進的波，據此判定\textbf{光本身就是一種電磁波}。

這是高中唯一系統介紹電磁學的章節。它解釋了發電、輸電、無線通訊、微波爐、X 光、手機等技術背後的同一套原理。馬克士威從方程式算出光速這件事，也直接埋下了\textbf{相對論}的開端。本章以\textbf{定性理解}為主。

\medskip
本章主線：

\begin{center}
電荷與電場 $\rightarrow$ 庫侖定律 $\rightarrow$ 電生磁（奧斯特）$\rightarrow$ 磁生電（電磁感應）$\rightarrow$ 馬克士威統一 $\rightarrow$ 電磁波與光 $\rightarrow$ 都卜勒效應
\end{center}

\section*{\sffamily 4-1\quad 電荷與電場}

\subsection*{\sffamily A.\ 電荷的三個基本性質}

把塑膠尺在頭髮上摩擦後能吸起小紙屑，這是最早被觀察到的電現象。實驗歸納出三個基本事實：

\begin{itemize}\setlength{\itemsep}{1pt}
\item \textbf{電荷有兩種}：稱為\textbf{正電荷}與\textbf{負電荷}（此命名出自富蘭克林，屬約定）。同號相斥、異號相吸。
\item \textbf{電荷守恆}：摩擦並非「製造」電荷，而是把電子從一個物體\textbf{轉移}到另一個物體，總電荷量不變。
\item \textbf{電荷量子化}：自然界的電荷都是基本電荷 $e\approx1.6\times10^{-19}$ C（庫侖）的整數倍。電子帶 $-e$，質子帶 $+e$。
\end{itemize}

原子核帶正電、電子帶負電，靠電力互相吸引而形成穩定的原子。這個「電力」就是本節要量化的庫侖力。

\subsection*{\sffamily B.\ 庫侖定律}

兩個帶電體若尺寸遠小於彼此距離，可視為\textbf{點電荷}。庫侖（Coulomb）以扭秤實驗測得：兩點電荷間的電力，與兩電荷量乘積成正比、與距離平方成反比。

\textbf{庫侖定律（Coulomb's law）}：真空中兩點電荷 $q_1$、$q_2$ 相距 $r$，彼此作用力大小為
$$\boxed{\;F = k\,\dfrac{q_1 q_2}{r^2}\;}$$
其中庫侖常數 $k\approx 9\times10^{9}$ N$\cdot$m$^2$/C$^2$。力的方向沿兩電荷連線：同號相斥、異號相吸。

\begin{補充框}{先補充：庫侖定律的使用要點}
\begin{itemize}\setlength{\itemsep}{1pt}
\item 公式裡的 $r$ 是\textbf{兩電荷之間的距離}，不是到原點或到某個面的距離。
\item 庫侖定律適用於\textbf{點電荷}或球對稱帶電球；不規則帶電體需用積分處理（高中不做）。
\item 計算時單位要先換成庫侖（C）與公尺（m）。例如 $1\,\mu$C$=10^{-6}$ C。
\end{itemize}
\textbf{注意：}庫侖定律與萬有引力 $F=G\,m_1m_2/r^2$ 形式相同，都是「平方反比」，但兩者強度差距極大：基本電力遠強於重力。日常之所以較少察覺強大的靜電力，是因為正、負電荷通常高度中和，對外幾乎抵消。
\end{補充框}

\subsection*{\sffamily C.\ 電場}

庫侖定律描述兩個電荷隔空互相施力，但未說明力如何「跨過」中間的空間傳遞。物理學以\textbf{場}的觀點來描述：電荷會在它周圍的空間建立一種狀態，稱為\textbf{電場（electric field）}。另一個電荷感受到的不是「遠方那顆電荷」，而是「它\textbf{所在位置}的電場」。

\textbf{電場的定義}：空間某點的電場 $\vec E$，定義為\textbf{單位正電荷在該點所受的電力}：
$$\vec E=\frac{\vec F}{q_0},\qquad \text{單位：N/C}.$$
由庫侖定律，點電荷 $Q$ 在距離 $r$ 處造出的電場大小為
$$\boxed{\;E=k\dfrac{|Q|}{r^2}\;}$$
方向：正電荷的電場向外輻射，負電荷的電場向內匯聚。一旦知道某點的 $\vec E$，任何電荷 $q$ 在該點受力即 $\vec F=q\vec E$。

\begin{補充框}{先補充：電場是什麼？「場」的概念}
\textbf{電場是「空間本身被電荷改變了狀態」的描述。} 在某點放一個試驗電荷，它會受力；這個「每單位電荷在該點所受的力」就是該點的電場。

引入場是為了取代\textbf{超距作用}的圖像。超距作用認為 A 直接、瞬間地對遠方的 B 施力，但這留下一個問題：力如何跨過中間的空間？場的觀點改為：A 先在它周圍\textbf{每一點}建立電場（改變了空間的狀態），B 只感受到它自己所在那一點的電場。力是「就地」作用的。（可類比為：你感受到的不是遠處的暖爐，而是身邊這塊空氣的溫度；暖爐先建立整個房間的「溫度場」，你只與身邊的空氣交互。）

\textbf{場是真實的物理實體，而非僅是計算工具。} 對靜止電荷，場確實可有可無——只用庫侖定律即可算力。但當電荷\textbf{運動}時，它造出的場不是瞬間改變遠方，而是以\textbf{光速}一波波傳出去；這個傳出去的場擾動可以完全脫離源頭，自己在真空中前進，這就是電磁波（見 4-5）。電磁波帶有能量與動量（陽光能曬熱皮膚），因此場不可能只是計算記號。本章末會回到這一點。
\end{補充框}

\subsection*{\sffamily D.\ 電力線}

法拉第（Faraday）以\textbf{電力線（electric field lines）}把電場視覺化，規則如下：

\begin{itemize}\setlength{\itemsep}{1pt}
\item 線上每一點的\textbf{切線方向}即該點的電場方向。
\item 線從\textbf{正電荷出發、向負電荷終止}（或延伸到無窮遠）。
\item \textbf{線越密的地方，電場越強}。
\item 電力線不會相交（否則交點會有兩個電場方向，互相矛盾）。
\end{itemize}

\fig{m4-fieldlines}{0.95}{圖 4-1：電力線。左——單一正電荷，電力線向外輻射；中——單一負電荷，電力線向內匯聚；右——一對正負電荷（電偶極），電力線由正電荷出、向負電荷入。}

\begin{補充框}{注意：電力線是視覺化工具，不是真實的「線」}
空間裡並沒有一條條實體的線。真實的電場是\textbf{每一點都有一個方向與大小的量}（向量場），布滿整個空間。電力線只是用來呈現場的方向與強弱（線越密代表越強）的畫法。
\end{補充框}

\section*{\sffamily 4-2\quad 電流的磁效應：電生磁}

\subsection*{\sffamily A.\ 奧斯特的發現}

在 1820 年以前，「電」與「磁」被視為兩門互不相干的學問。1820 年，丹麥物理學家\textbf{奧斯特（Oersted）}在演示電流時注意到：旁邊的指南針偏轉了。電路一通電，磁針就偏轉；一斷電，磁針回正。這顯示出一個結論：

\begin{center}
\textbf{電流會產生磁場。}
\end{center}

這是電與磁有關聯的第一項直接證據，也是電磁學統一的起點。

\subsection*{\sffamily B.\ 直導線的磁場與右手定則}

實驗顯示，一條通有電流的直導線，周圍的磁場是\textbf{一圈圈以導線為中心的同心圓}。磁場方向以\textbf{安培右手定則}判定。

\begin{補充框}{先補充：（電生磁的）右手定則怎麼用}
\textbf{右手大拇指指向電流方向，其餘四指自然彎曲環繞的方向，就是磁場的方向。}

由此可知：直導線周圍的磁場\textbf{不是直線，而是環繞導線的同心圓}。磁場大小隨距離增加而減弱（離導線越遠越弱），也隨電流增大而增強。\textbf{注意：}本章一律使用右手定則（拇指電流、四指磁場），不要與左手定則混用。
\end{補充框}

\fig{m4-oersted}{0.92}{圖 4-2：左——電流由紙面射出（中心圓點）時，周圍磁場 $B$ 為同心圓；右——安培右手定則：右手握住導線，拇指指電流方向，四指環繞方向即磁場方向。}

把導線繞成\textbf{線圈（螺線管）}，各圈的磁場疊加，內部即形成一個很強、近乎均勻的磁場，這就是\textbf{電磁鐵}：通電有磁、斷電無磁，磁力強弱可用電流控制。電鈴、繼電器、磁浮列車、廢鐵場的吊磁都應用此原理。

\begin{補充框}{注意：磁場沒有「磁單極」}
找不到單獨存在的 N 極或 S 極。把磁鐵切成兩半，每一半仍各有 N、S 兩極。這一點在 4-4 會再用到。
\end{補充框}

\section*{\sffamily 4-3\quad 電磁感應：磁生電}

電能生磁，反過來磁能不能生電？這是 1820 年代之後物理學家致力研究的問題，最後由\textbf{法拉第（Faraday）}找到答案。

\subsection*{\sffamily A.\ 關鍵在「變動」}

法拉第發現：把磁鐵\textbf{插入或抽出}線圈時，線圈會短暫產生電流（\textbf{感應電流}）；但磁鐵若\textbf{靜止不動}地放在線圈裡，無論磁多強，都沒有電流。

\textbf{電磁感應（electromagnetic induction）}：當穿過線圈的磁場（磁力線數）\textbf{發生變動}時，線圈中會產生\textbf{感應電動勢}，驅動\textbf{感應電流}。關鍵在「變動」：磁場變化越快，感應電流越大；磁場不變，就沒有感應電流。

換句話說，\textbf{變動的磁場會在周圍產生電場}，這個電場推動了導線裡的電荷。

\fig{m4-induction}{0.62}{圖 4-3：電磁感應。磁鐵相對線圈運動（如插入）時，穿過線圈的磁場發生變動，線圈中產生感應電流（電流計 G 偏轉）。磁鐵靜止時則無感應電流。}

\subsection*{\sffamily B.\ 電磁感應的應用}

電磁感應是人類取得電力的根本方式。火力、水力、核能、風力等各種發電做的都是同一件事：\textbf{設法讓線圈與磁場之間持續相對運動}，把機械能轉換成電能。

\begin{center}
\renewcommand{\arraystretch}{1.35}
\begin{tabular}{p{0.20\linewidth} p{0.46\linewidth} p{0.22\linewidth}}
\hline
\textbf{裝置} & \textbf{原理} & \textbf{能量轉換} \\
\hline
發電機 & 線圈在磁場中轉動，穿過線圈的磁場不斷變動，持續產生感應電流 & 機械能 $\rightarrow$ 電能 \\
變壓器 & 初級線圈的變動電流產生變動磁場，在次級線圈感應出電壓 & 改變交流電壓 \\
感應爐、無線充電 & 變動磁場在金屬或線圈中感應電流 & 隔空傳能 \\
\hline
\end{tabular}
\end{center}

\begin{補充框}{注意：發電機是能量轉換，不是無中生有}
磁鐵靜止不動時不會發電，必須有「變動」。發電機並非憑空造出能量，而是\textbf{轉換}能量：需要外力去轉動它（這也是腳踏車發電時感到「變重」的原因）。整個過程遵守能量守恆。
\end{補充框}

\section*{\sffamily 4-4\quad 馬克士威的統一}

到 1860 年代，零散的電磁實驗已累積成一批規律：庫侖定律、奧斯特（電生磁）、法拉第（變動磁場生電）等。\textbf{馬克士威（Maxwell）}把它們整理成\textbf{一組四條方程式}，並補上關鍵的最後一項。

\subsection*{\sffamily A.\ 補上的一項：變動電場也產生磁場}

法拉第指出「變動磁場產生電場」。馬克士威從數學的對稱與一致性推出：必定還有一條對稱的規律——

\begin{center}
\textbf{變動的電場，也會產生磁場。}
\end{center}

這條規律在當時尚無直接實驗、純由理論推出，是整個統一的樞紐。於是電與磁的關係變得完全對稱：
$$\text{變動磁場}\ \longrightarrow\ \text{電場};\qquad \text{變動電場}\ \longrightarrow\ \text{磁場}.$$

\subsection*{\sffamily B.\ 馬克士威方程式的定性意義}

四條方程式不需背公式，但要理解它們各自的意義：

\begin{itemize}\setlength{\itemsep}{2pt}
\item \textbf{電荷產生電場}（高斯定律）：電力線從正電荷出、向負電荷入，即庫侖定律的場版本。
\item \textbf{沒有磁單極}：磁力線永遠是封閉迴圈，找不到單獨的 N 極或 S 極。
\item \textbf{變動磁場產生電場}（法拉第定律）：這是電磁感應、發電機的根據。
\item \textbf{電流與變動電場產生磁場}（安培-馬克士威定律）：奧斯特的電生磁，再加上馬克士威補的一項。
\end{itemize}

這四條合起來涵蓋了所有古典電磁現象，從雷電到原子內電子所受的力皆包含在內。馬克士威幾乎沒有做新實驗，主要貢獻是\textbf{理論綜合}：把既有的實驗規律統整、補全，並推出全新的預測（下一節的電磁波）。這展現了把「看似無關的兩類現象，證明其實是同一件事的兩面」的研究路徑，也是物理學追求「統一」的範例。

\section*{\sffamily 4-5\quad 電磁波與光}

\subsection*{\sffamily A.\ 電磁波如何傳播}

把 4-4 的兩條規律接在一起：

\begin{center}
變動電場 $\rightarrow$ 變動磁場 $\rightarrow$ 變動磁場再產生變動電場 $\rightarrow$ 又產生磁場 $\rightarrow$ $\cdots$
\end{center}

電場與磁場互相激發、彼此接力，一旦啟動，這個擾動就能\textbf{自行往外傳播、脫離源頭、在空間中前進}，這就是\textbf{電磁波（electromagnetic wave）}。

\textbf{電磁波的定義}：電磁波是\textbf{電場 $E$ 與磁場 $B$ 互相垂直、且都垂直於傳播方向}的橫波：
$$E\perp B\perp \text{傳播方向}.$$
它\textbf{不需要任何介質}，可以在真空中傳播。其在真空中的速率對所有電磁波都相同，記為 $c$。

\fig{m4-emwave}{0.66}{圖 4-4：電磁波是橫波。電場 $E$（紅）與磁場 $B$（藍）互相垂直，且都垂直於傳播方向 $x$；兩者同步地週期起伏，向前傳播。}

\begin{補充框}{先補充：電磁波不需要介質，那它在「振動」什麼？}
聲波靠空氣分子振動傳遞、水波靠水傳遞，但電磁波可以穿越真空（陽光即是穿過太空真空到達地球）。在電磁波中\textbf{振動的不是任何物質，而是電場與磁場本身}。

在空間某一點，\textbf{電場的大小隨時間上下擺動}（如正弦曲線），它一變動就在旁邊激發出同樣在擺動的磁場，磁場的變動又往前激發新的電場。因此「在振動的」是這一點的\textbf{電場強度與磁場強度}在週期變化，而非任何粒子在移動。真空沒有物質，但容得下電磁場；場本身就能起伏並向前傳遞。

\textbf{歷史上「以太」的假設已被推翻。} 十九世紀的物理學家假設宇宙充滿一種看不見的介質「以太（ether）」，認為光是以太的振動。但麥克生-莫立實驗（Michelson-Morley）想測出地球相對以太運動造成的差異，結果\textbf{量不到任何效應}。結論是：\textbf{並沒有以太這種介質}，電磁波振動的就是場本身，不需要載體。\textbf{注意：}不能由「以太不存在」推論「場也是假的」——場是真實的，只是它\textbf{不需要}一個物質載體。
\end{補充框}

\subsection*{\sffamily B.\ 波速、頻率、波長的關係}

電磁波是波，遵守一切週期波的共通關係。

\begin{補充框}{先補充：波速、頻率、波長的關係}
對任何週期波，
$$\boxed{\;v=f\lambda\;}$$
其中 $v$ 是波速，$f$ 是\textbf{頻率（frequency）}（每秒振動次數，單位 Hz），$\lambda$ 是\textbf{波長（wavelength）}（一個完整波的長度）。

對真空中的電磁波，$v=c\approx 3\times10^{8}$ m/s（光速），所以 $c=f\lambda$。由於 $c$ 固定，\textbf{頻率越高，波長越短}；反之亦然。
\end{補充框}

\subsection*{\sffamily C.\ 光就是電磁波}

馬克士威從方程式推導電磁波時\textbf{算出了它的速率}，而這個速率\textbf{只由兩個電學、磁學常數決定}，算出的結果恰好等於 $3\times10^8$ m/s。當時這個數值早已被獨立測出，是\textbf{光速}。

兩個數值相同並非巧合。馬克士威據此推論：

\begin{center}
\textbf{光，就是一種電磁波。}
\end{center}

其後\textbf{赫茲（Hertz）}在實驗室中產生並接收到電磁波，證實了這個推論。從此「光學」併入「電磁學」，是又一次統一。

\begin{延伸框}{馬克士威如何從方程式得到光速}
這部分是已確立的科學史，可據實了解其推導。

電磁學裡有兩個從\textbf{靜態實驗}就能測出的基本常數，原本與光無關：
\begin{itemize}\setlength{\itemsep}{1pt}
\item \textbf{真空電容率 $\varepsilon_0$}：描述電場在真空中有多容易建立（出現在庫侖定律，$k=1/4\pi\varepsilon_0$）。它來自「電」的實驗（量電荷間的力）。
\item \textbf{真空磁導率 $\mu_0$}：描述電流產生磁場有多容易（出現在電流磁效應的公式中）。它來自「磁」的實驗（量電流間的力）。
\end{itemize}
這兩個常數都沒有用到任何光學儀器。馬克士威把四條方程式組合起來，數學上會得到一個\textbf{波動方程式}；任何波動方程式都自帶一個波速，而這裡算出的波速恰好是
$$\boxed{\;c=\frac{1}{\sqrt{\varepsilon_0\,\mu_0}}\;}$$
代入當時測得的 $\varepsilon_0$、$\mu_0$，得到約 $3\times10^8$ m/s，與已測得的光速一致。重點不必背公式或常數數值，只需掌握一句話：\textbf{「電的常數」與「磁的常數」相乘、開根號、取倒數，就得到光速}。電與磁的內在性質裡，本來就決定了電磁擾動傳播的速率，而那正是光速。
\end{延伸框}

\begin{延伸框}{為什麼這個推論重要：相對論的開端}
這同樣是已確立的科學史。馬克士威算出的 $c=1/\sqrt{\varepsilon_0\mu_0}$ 是一個\textbf{固定的數}，公式裡沒有任何「相對於誰」的字眼。

在馬克士威之前的力學中，速度永遠是相對的：在時速 100 的火車上朝前以時速 10 丟球，地面上量到的球速是 110，速度會隨觀察者的運動疊加（伽利略速度相加）。但 $c=1/\sqrt{\varepsilon_0\mu_0}$ 只給一個值，未說明相對於哪個參考系。這帶來一個矛盾：若光速也會疊加，則「追著光跑」的人應量到較慢的光、「迎向光」的人量到較快的光，那 $c$ 就不該是固定值。當時的補救是假設有以太作為絕對參考系，但麥克生-莫立實驗測不到以太。

物理學因此面臨抉擇：放棄極為成功的馬克士威方程式，或放棄「速度可任意疊加」這個自牛頓以來的假設。1905 年，愛因斯坦（Einstein）選擇保留馬克士威，並把「光速對所有慣性觀察者都是同一個 $c$」立為公設（光速不變原理），由此推出狹義相對論：時間膨脹、長度收縮、同時性是相對的。因此，馬克士威算出一個「不說相對於誰」的固定光速，正是促成相對論的起點。今天的 SI 單位制甚至反過來把光速 $c$ 定義為固定值（$299\,792\,458$ m/s），用它來定義「公尺」。
\end{延伸框}

\subsection*{\sffamily D.\ 電磁波譜}

可見光只是電磁波家族中很窄的一段。所有電磁波\textbf{本質相同}（都是 $E$、$B$ 互生的橫波、都以光速行進），差別只在\textbf{頻率（波長）不同}。

\fig{m4-spectrum}{0.98}{圖 4-5：電磁波譜（橫軸為波長，對數刻度）。由左到右波長漸短、頻率漸高、光子能量漸大：無線電波 $\rightarrow$ 微波 $\rightarrow$ 紅外線 $\rightarrow$ 可見光 $\rightarrow$ 紫外線 $\rightarrow$ X 射線 $\rightarrow$ $\gamma$ 射線。可見光只占其中極窄的一段。}

\begin{center}
\renewcommand{\arraystretch}{1.3}
\begin{tabular}{p{0.20\linewidth} p{0.28\linewidth} p{0.40\linewidth}}
\hline
\textbf{波段} & \textbf{波長範圍（約）} & \textbf{日常應用} \\
\hline
無線電波 & 公尺以上 & 廣播、電視、手機通訊 \\
微波 & 公分至公尺 & 微波爐、Wi-Fi、雷達、衛星 \\
紅外線 & $\mu$m 級 & 熱感應、夜視、遙控器 \\
可見光 & $400\sim700$ nm & 人眼可見、照明、光學 \\
紫外線 & nm 級 & 殺菌、驗鈔 \\
X 射線 & $0.01\sim10$ nm & 醫學造影、安檢 \\
$\gamma$ 射線 & $<0.01$ nm & 放射治療、核醫 \\
\hline
\end{tabular}
\end{center}

越往高頻，光子能量越大，穿透力與破壞力也越強，因此須注意 X 光、$\gamma$ 射線的劑量。

\begin{補充框}{注意：關於電磁波譜的幾點}
\begin{itemize}\setlength{\itemsep}{1pt}
\item 無線電波與 X 光\textbf{本質完全相同}，只是頻率不同的電磁波。
\item 可見光只占電磁波譜中極窄的一段，人眼可見的範圍很有限。
\item 微波爐並非用核輻射加熱。微波是\textbf{非游離輻射}，能量遠低於 X、$\gamma$ 射線，僅讓水分子加速振動而生熱。
\end{itemize}
\end{補充框}

\section*{\sffamily 4-6\quad 都卜勒效應}

\subsection*{\sffamily A.\ 現象}

救護車朝你駛來時鳴笛聽起來較\textbf{高（尖）}，經過後遠離時較\textbf{低（沉）}。車上鳴笛的頻率並未改變，但你聽到的頻率改變了。這就是\textbf{都卜勒效應（Doppler effect）}。

\textbf{都卜勒效應（定性）}：當波源與觀察者有\textbf{相對運動}時，觀察者測到的頻率會偏離波源本身的頻率：

\begin{itemize}\setlength{\itemsep}{1pt}
\item \textbf{互相靠近} $\rightarrow$ 波前被壓縮 $\rightarrow$ 波長變短 $\rightarrow$ 頻率變高。
\item \textbf{互相遠離} $\rightarrow$ 波前被拉伸 $\rightarrow$ 波長變長 $\rightarrow$ 頻率變低。
\end{itemize}

\fig{m4-doppler}{0.66}{圖 4-6：移動的波源。波源一邊發波一邊前進，使前方的波前一個個擠近（波長短、頻率高，藍移），後方的波前被拉開（波長長、頻率低，紅移）。}

\begin{補充框}{先補充：都卜勒效應的成因}
鳴笛器每秒發出固定數目的聲波（一圈圈同心圓波前）。若波源不動，各圈波前都從同一中心向外擴，四面八方波長相同。若波源往前移動，每發出新的一圈時，波源已往前挪了一段，使新波前的中心跟著前移：

\begin{itemize}\setlength{\itemsep}{1pt}
\item \textbf{波源前方}：波前的中心一直往前追，波前彼此擠近，波長變短、頻率變高。
\item \textbf{波源後方}：波前的中心一直往後離開，波前彼此拉開，波長變長、頻率變低。
\end{itemize}
（可類比為：有人朝你走來並每秒丟一顆球，因為他越走越近，球到手的間隔越來越短，你接球的頻率變高；他走過後遠離，間隔被拉長，接球變慢。聲波的波前就相當於那一顆顆球。）

\textbf{注意：}
\begin{itemize}\setlength{\itemsep}{1pt}
\item 改變的是\textbf{觀察者接收到的頻率}，波源本身發出的頻率並未改變（駕駛自己聽到的笛聲始終一樣）。
\item 音調是在「波源通過你的那一刻」前後，由偏高切換為偏低，而非漸漸變低；靠近時一直偏高、遠離時一直偏低。
\item 音調（高低）由都卜勒效應造成；音量（大小聲）則由距離遠近造成，兩者是不同的事。
\end{itemize}
\end{補充框}

\subsection*{\sffamily B.\ 光的都卜勒效應：紅移與藍移}

聲波會發生都卜勒效應，光（電磁波）也會。當光源與觀察者相對運動：

\begin{itemize}\setlength{\itemsep}{1pt}
\item \textbf{遠離} $\rightarrow$ 波長變長 $\rightarrow$ 光譜往紅端偏移 $\rightarrow$ 稱\textbf{紅移（redshift）}。
\item \textbf{靠近} $\rightarrow$ 波長變短 $\rightarrow$ 光譜往藍端偏移 $\rightarrow$ 稱\textbf{藍移（blueshift）}。
\end{itemize}

光的都卜勒效應有兩個重要應用：

\begin{itemize}\setlength{\itemsep}{1pt}
\item \textbf{測速雷達}：朝車輛發射電磁波並接收反射波。迎面駛來的車使反射波頻率升高（藍移），遠離的車使反射波頻率降低（紅移）。量出反射波與發射波的頻率差，即可換算出車速。
\item \textbf{宇宙膨脹}：天文學家觀測遙遠星系，發現其光譜幾乎都呈\textbf{紅移}，且越遠的星系紅移越大。這表示這些星系正在遠離我們、且越遠者遠離越快，即哈伯（Hubble）發現的\textbf{宇宙正在膨脹}，是大霹靂理論的重要觀測支柱。
\end{itemize}

\begin{補充框}{注意：「紅移」的意思}
「紅移」不是指「光真的變成紅色」，而是\textbf{整個光譜往長波長（低頻）方向移動}。可見光的紅移有可能移出可見範圍而成為紅外線。
\end{補充框}

\section*{\sffamily 4-7\quad 本章重點整理}

\begin{補充框}{一頁複習（公式與關鍵結論）}
\begin{itemize}\setlength{\itemsep}{2pt}
\item \textbf{電荷}：兩種（正/負）、守恆、量子化（基本電荷 $e\approx1.6\times10^{-19}$ C）；同號斥、異號吸。
\item \textbf{庫侖定律}：$F=k\,q_1q_2/r^2$，$k\approx9\times10^9$ N$\cdot$m$^2$/C$^2$；平方反比。
\item \textbf{電場}：$\vec E=\vec F/q_0$，點電荷 $E=k|Q|/r^2$；電力線正出負入、越密越強。場是真實的物理實體。
\item \textbf{電生磁}（奧斯特）：直導線周圍為環形磁場，以右手定則判定（拇指電流、四指磁場）；線圈即電磁鐵。
\item \textbf{磁生電}（法拉第，電磁感應）：磁場\textbf{變動}才產生感應電流；發電機、變壓器的根據。
\item \textbf{馬克士威統一}：補上「變動電場生磁場」，四條方程式涵蓋所有古典電磁現象。
\item \textbf{電磁波}：$E\perp B\perp$ 傳播方向的橫波，不需介質，真空中速率為 $c\approx3\times10^8$ m/s；$v=f\lambda$（即 $c=f\lambda$）。
\item \textbf{光是電磁波}：馬克士威算出波速 $=$ 光速（$c=1/\sqrt{\varepsilon_0\mu_0}$），並埋下相對論的開端。電磁波譜：無線電 $\rightarrow$ 微波 $\rightarrow$ 紅外 $\rightarrow$ 可見 $\rightarrow$ 紫外 $\rightarrow$ X $\rightarrow$ $\gamma$，頻率漸高、能量漸大。
\item \textbf{都卜勒效應}：相對靠近 $\rightarrow$ 頻率變高（藍移），遠離 $\rightarrow$ 變低（紅移）；應用於測速雷達、宇宙膨脹。
\end{itemize}
\end{補充框}

\end{document}
